\documentstyle[12pt]{article}
\oddsidemargin 0.1in
\evensidemargin -0.8in
\topmargin -0.75in
\textheight 24cm
\textwidth 16cm


\def\Vec#1{\mbox{$\bf#1$}}
\newcommand{\Sum}{\displaystyle\sum}
\newcommand{\Eqref}[1]{(\ref{#1})}
\newcommand{\Frac}[2]{\mbox{$\displaystyle\frac{#1}{#2}$}}


\title{\bf Analytical Approach for Simplifying\\
           Dynamical Systems of Polynomial Type}

 \author{
{\small
 \begin{tabular}{c}
    J. Palaci\'an and P. Yanguas
 \end{tabular}} \\[2ex]
{\footnotesize
 \begin{tabular}{c}
               {\it Departamento de Matem\'atica e Inform\'atica,}\\
               {\it Universidad P\'ublica de Navarra, 31006 Pamplona, Spain}
   \end{tabular}}}
\date{}



\begin{document}
\maketitle

Many dynamical systems are modelled by a system of ordinary differential
equations which is formed by the sum of a linear part plus a small
perturbation.
Moreover, the perturbation is a finite sum of nonlinear vector fields which
begins
with quadratic polynomials. In other words, we study systems of the type
\begin{equation}
\Frac{d\,\Vec{x}(t)}{d\,t}\,=\,A\,\Vec{x}(t)+
   \Sum_{n=1}^L\Frac{\varepsilon^n}{n!}\Vec{F}_n(\Vec{x}(t)),\label{sys}
\end{equation}
where $A$ is an $m\times m$--matrix with constant coefficients with physical
dimensions [time$^{-1}$], $\varepsilon$ stands for a dimensionless small
parameter
and the vector field $\Vec{F}_n$ has $m$ components and each component
is a polynomial of degree $n+1$.

Analytical methods to deal with systems of the type~\Eqref{sys} are based
on the
fact  that these differential equations can be understood as perturbations
of the principal part, $A\,\Vec{x}(t)$. In such circumstances, one can think of
reducing the problem to a simpler one by means of adequate changes of
variables.

These equations appear in many situations of Dynamical Systems Theory, such as
stability and bifurcation analysis of equilibrium points and periodic orbits or
singularity theories. Besides, applications are encountered in several
fields such
as Engineering (van der Pol--like and Duffing--like equations, see~[1]) or
Astrodynamics
(potential functions modelling elliptical galaxies, see~[7]).

The first step in the analysis of Equation~(\ref{sys}) was given by
Poincar\'e in~[6]. He considered the problem of reducing it to a system of
linear equations
$$\Frac{d\,\Vec{y}(t)}{d\,t}\,=\,A\,\Vec{y}(t),$$
by defining the formal change of variables
$\Vec{y}(t)=\Vec{\varphi}(\Vec{x}(t);\varepsilon)=\Vec{x}(t)+\ldots$.
Concretely, he
found a solution in the case where $A$ were diagonalizable and its eigenvalues
$\lambda_1,\ldots,\lambda_m$ satisfied the resonant condition, that is,
$\lambda_j\neq\Sum
k_i\,\lambda_i$ for
$j=1,\ldots,m$ and all integer vectors $k=(k_1,\ldots,k_m)$ with $k_i\ge0$ and
$|k|=k_1+\ldots+k_m\ge 2$. He also proved that, if in addition to
the above, the eigenvalues of $A$, considered in the complex plane, lie
strictly to one side
of a line through the origin of the plane, then the formal series  which
gives the change of variables converges (for a proof see reference~[7]).

When these conditions for the eigenvalues are not satisfied, one tries to
find the
normal form of~\Eqref{sys}, that is, to reduce this system to an
``equivalent" one
\begin{equation}
\Frac{d\,\Vec{y}(t)}{d\,t}\,=\,A\,\Vec{y}(t)+
   \Sum_{n=1}^M\Frac{\varepsilon^n}{n!}\Vec{G}_n(\Vec{y}(t))+{\cal
O}(\varepsilon^{M+1}),\label{sys2}
\end{equation}
where $M\ge L$ and $\Vec{G}_n$ are ``simpler" than $\Vec{F}_n$ for
$n=1,\ldots,L$.

Meyer~[3] has generalized Poincar\'e's results by giving the Normal Form
Theorem, based on Lie
transformations. See reference~[2] for the method applied to Hamiltonian
systems and
reference~[4] for the generalization to differential equations and tensor
fields. We summarize THE
theorem briefly. Let $A=A_S+A_N$ be the Jordan decomposition of a square
matrix into its semisimple
and nilpotent parts, respectively. Let $A^t$ be the transpose matrix of $A$
(which is, indeed, the
adjoint matrix for the standard inner product
in ${\bf R}^m$). Denote by $[\,\cdot\,,\,\cdot\,]$ the Lie bracket of two
vector fields, that is,
$[\,\Vec{g}_1\,,\,\Vec{g}_2\,]=(\partial{\Vec{g}_1}/\partial{\Vec{x}})\,\Vec{g}_
2
-(\partial{\Vec{g}_2}/\partial{\Vec{x}})\,\Vec{g}_1$
and by
${\cal L}_A$ the linear operator associated to $A$, such that, for any
vector field
$\Vec{g}$ defined in a domain of
${\bf R}^m$,
${\cal L}_A(\Vec{g})=[\,\Vec{g}\,,\,A\,\Vec{x}\,]$. The theorem states that
there exists a formal transformation
$\Vec{y}(t)=\Vec{\varphi}(\Vec{x}(t);\varepsilon)=\Vec{x}(t)+\ldots$ which
transforms~\Eqref{sys} into~\Eqref{sys2}, where each term $\Vec{G}_n$, $n=1,
\ldots,M$  belongs to the kernel of the operator ${\cal L}_{A^t}$, that is,
\begin{equation}\label{theo}
[\,\Vec{G}_n\,,\,A^t\,\Vec{x}\,]=0, \quad n=1,\ldots,M.\end{equation}

Note that the transformation is analytic since we truncate the change of
variables at an order $M$. Moreover, by virtue
of~(\ref{theo}), it is not difficult to see that,
$[\,A\,\Vec{x}\,,\,A_S\,\Vec{x}\,]=0$ and $[\,\Vec{G}_n\,,\,A_S\,\Vec{x}\,]=0$,
$n=1,\ldots, M$, since $\mbox{ker}\,{\cal L}_{A^t}\subset\mbox{ker}\,{\cal
L}_{A_S}$.

However, if $A_S=0$, the consequence of applying the Normal Form Theorem is
not so drastic.
Then, the strategy to reduce the system must be different. Moreover, even
if $A_S\neq0$, the
above result does not permit to further simplify the initial
system. Therefore, one can think of generalizing the calculation of normal
forms by Lie transformations to simplify  systems of the type~(\ref{sys})
even more.

Let $T$ be a square $m\times m$--matrix with Jordan decomposition $T=T_S+T_N$,
that is, $T_S$ semisimple and $T_N$ nilpotent. In each step of the process of
a Lie transformation ($n=1,\ldots,M$) one has to solve the homology equation
$${\cal L}_A(\Vec{W}_n)+\Vec{G}_n\,=\,\widetilde{\Vec{F}}_n,$$
where $\Vec{W}_n$ is related to the $n$--th order of the generating
function whereas
$\widetilde{\Vec{F}}_n$ denotes the vectors computed in the previous steps.
Then, we try to find $\Vec{G}_n$ such that $[\,\Vec{G}_n\,,\,T\,\Vec{x}\,]=0$.
For that,
we split
$\widetilde{\Vec{F}}_n=\widetilde{\Vec{F}}^{\#}_n+\widetilde{\Vec{F}}^{\&}_n$ wh
ere
$\widetilde{\Vec{F}}^{\#}_n\in
\mbox{ker}\,({\cal L}_T)$, i.e., ${\cal
L}_T(\widetilde{\Vec{F}}^{\#}_n)=[\,\widetilde{\Vec{F}}^{\#}_n\,,\,T\,\Vec{x}\,]
=0$, identify
$\Vec{G}_n=\widetilde{\Vec{F}}^{\#}_n$ and
$\widetilde{\Vec{F}}^{\&}_n=\widetilde{\Vec{F}}_n-\widetilde{\Vec{F}}^{\#}_n$. T
he vector
$\Vec{W}_n$ must be a solution of the system of partial differential equations
${\cal
L}_A(\Vec{W}_n)=\widetilde{\Vec{F}}^{\&}_n.$

If, in addition to the above,  $T$ has been chosen so that it commutes with
$A$, then
$[\,A\,\Vec{x}\,,\,T\,\Vec{x}\,]=0$ and one reduces the original system.
Note that this
approach enlarges the Normal Form Theorem and generalizes it since one does not
need to take $A^t$ but any other matrix $T$.

One can prove that if $T$ satisfies: i) it is semisimple, that is, $T_N=0$;
ii) $A\,T=T\,A$  and
iii) $\mbox{ker}\,({\cal L}_A)\cap\mbox{im}\,({\cal L}_T)=\{0\}$, then the
decomposition of
$\widetilde{\Vec{F}}_n$ is performed straightforwardly. Moreover the
 system of partial differential  equations is solvable and $\Vec{W}_n$ are
polynomial
vector fields, for details see~[6]. In this ``optimal"  case one can reduce the
initial system to an equivalent one, applying several  Lie transformations.
Now, the final reduced
system can be studied more easily and the reconstruction of the original
equations is performed by
making use of the inverse Lie transformation up to a certain order
$M$.

If some condition of the latter paragraph does not hold, one must
apply the Lie transformations carefully. Now, the process can terminates
before,
either because it is not possible to find out more semisimple matrices,
$T$, or, because
in any Lie transformation, one can compute $\Vec{W}_n$ only up to a certain
order
$M$ so that the generating functions are polynomials.

We illustrate the technique with some examples of differential systems
putting special emphasis in the implementation of the method with the
algebraic
processor {\sc Mathematica}.



 {\small
\begin{thebibliography}{}

\item[1.]~T.~de Zeeuw, Motion in the Core of a Triaxial Potential, {\sl
Monthly
Notices of the Royal Astronomical Society\/}
{\bf 215}, 731--760 (1985).

\item[2.]~A.~Deprit, Canonical Transformations Depending on a Small Parameter,
{\sl Celestial Mechanics\/} {\bf 1}, 12--30
(1969).

\item[3.]~K.~Meyer, Normal Forms for the General Equilibrium,
{\sl Funkcialaj Ekvacioj\/} {\bf 27}, 261--271 (1984).

\item[4.]~K.~Meyer, A Lie Transform Tutorial~II,
in {\sl Computer Aided Proofs in Analysis}, eds. K.~Meyer and
D.~S.~Schmidt,  The IMA Volumes in Mathematics and its Applications {\bf 28},
Springer--Verlag, New York, 190--210 (1991).

\item[5.]~J.~Palaci\'an and P.~Yanguas, Reduction of Polynomial Differential
Equations by Constructing Formal Integrals, in preparation.

\item[6.]~H.~Poincar\'e, Memoire sur les Courbes Definie por une Equation
Differentialle IV, {\sl J. Math. Pures Appl.\/} {\bf 1}, 167--244 (1885).

\item[7.]~J.~A.~Sanders and F.~Verhulst, {\sl Averaging Methods in Nonlinear
Dynamical Systems}, Applied Mathematical Sciences {\bf 59},
Springer--Verlag, New
York (1985).
\end{thebibliography}}

\end{document}

